\documentclass{article}

% Recommended, but optional, packages for figures and better typesetting:
\usepackage{microtype}
\usepackage{graphicx}
\usepackage{subcaption}
\usepackage{booktabs} % for professional tables

% hyperref makes hyperlinks in the resulting PDF.
\usepackage{hyperref}

% Attempt to make hyperref and algorithmic work together better:
\newcommand{\theHalgorithm}{\arabic{algorithm}}

% Use the following line for the initial blind version submitted for review:
\usepackage{icml2026}
\renewcommand{\ttdefault}{lmtt} % Use Latin Modern Typewriter to avoid Courier font metric issues

\usepackage{amsmath}
\usepackage{amssymb}
\usepackage{mathtools}
\usepackage{amsthm}

% if you use cleveref..
\usepackage[capitalize,noabbrev]{cleveref}

%%%%%%%%%%%%%%%%%%%%%%%%%%%%%%%%
% THEOREMS
%%%%%%%%%%%%%%%%%%%%%%%%%%%%%%%%
\theoremstyle{plain}
\newtheorem{theorem}{Theorem}[section]
\newtheorem{proposition}[theorem]{Proposition}
\newtheorem{lemma}[theorem]{Lemma}
\newtheorem{corollary}[theorem]{Corollary}
\theoremstyle{definition}
\newtheorem{definition}[theorem]{Definition}
\newtheorem{assumption}[theorem]{Assumption}
\theoremstyle{remark}
\newtheorem{remark}[theorem]{Remark}

% The \icmltitle you define below is probably too long as a header.
% Therefore, a short form for the running title is supplied here:
\icmltitlerunning{\smash{Spectral Curvature Alignment}}

\begin{document}

\twocolumn[
  \icmltitle{Spectral Curvature Alignment: Resolving Anisotropic Representation\texorpdfstring{\\}{ }
    Interference in Multi-Task Model Merging}

  \begin{icmlauthorlist}
    \icmlauthor{Anonymous Author(s)}{anon}
  \end{icmlauthorlist}

  \icmlaffiliation{anon}{Anonymous Institution}
  \icmlcorrespondingauthor{Anonymous Author}{anon@example.com}
  \icmlkeywords{Model Merging, Spectral Methods, Curvature Alignment, representation collapse}

  \vskip 0.3in
  \begin{abstract}
    Multi-task model merging has emerged as a powerful paradigm for combining multiple specialized expert models into a single multi-task network without costly retraining. 
    However, standard weight averaging (WA) often suffers from severe performance degradation due to destructive interference among task-specific representations. 
    While recent heuristics, such as Isotropic Parameter Resonance (IPR), attempt to heal this ``representation collapse'' by uniformly scaling task updates, they rest on the highly restrictive assumption that the local loss landscape is isotropic.
    In this work, we propose \textbf{Spectral Curvature Alignment (SCA)}, a mathematically rigorous framework that relaxes this assumption by incorporating second-order curvature information. 
    By modeling the multi-task loss surface quadratically using the diagonal Fisher Information Matrix, we solve for the mathematically optimal, layer-wise spectral scale factors of the merged update's singular value decomposition (SVD). 
    We identify a fundamental challenge: cross-task destructive interference systematically under-scales the global update magnitude when optimizing local curvature directly. 
    To resolve this, we formulate \textbf{SCA-IPR}, which decouples directional spectral alignment (via SCA) from global energy calibration (via isotropic resonance).
    We prove that the resulting regularized solver acts as a biased estimator that minimizes estimation variance under noisy Fisher measurements.
    Empirically, on a ResNet-18 benchmark across MNIST, FashionMNIST, and CIFAR-10, SCA-IPR significantly outperforms WA by \textbf{+20.37\%} and the state-of-the-art U-IPR baseline by \textbf{+0.45\%}, demonstrating the theoretical and practical necessity of anisotropic curvature alignment in model fusion.
  \end{abstract}
]

\printAffiliationsAndNotice{} % this must go after the closing bracket ] of \twocolumn[

\section{Introduction}
\label{submission:intro}

Deep learning has achieved remarkable success through fine-tuning pre-trained foundation models—such as massive Transformers \cite{vaswani2017attention, devlin2018bert, brown2020language}, early sequence-to-sequence or attention models \cite{sutskever2014sequence, bahdanau2014neural}, or deep residual networks \cite{he2016deep, radford2021learning}—on specialized downstream tasks. 
To adapt these models efficiently, parameter-efficient fine-tuning (PEFT) approaches such as LoRA \cite{hu2021lora}, Prefix-Tuning \cite{li2021prefix}, and Adapters \cite{houlsby2019parameter} are widely used to construct task-specific vectors \cite{ilharco2022editing}. 
While this paradigm yields highly capable expert models, deploying separate parameters for each task incurs massive storage and computational overhead, violating standard scaling laws \cite{kaplan2020scaling} and limiting multi-task scalability \cite{caruana1997multitask, shazeer2017outrageously}. 
Multi-task model merging offers an elegant alternative: it fuses several fine-tuned expert models sharing a common progenitor (or initialization) back into a single multi-task model, entirely bypassing the need for multi-task training or joint fine-tuning \cite{wortsman2022model, ainsworth2023git}.

Historically, standard Weight Averaging (WA) serves as the primary baseline for model merging \cite{izmailov2018averaging, wortsman2022model}. 
While WA succeeds when the specialized experts reside within the same basin of the loss landscape, it breaks down catastrophically when merging models fine-tuned on highly dissimilar or disjoint tasks \cite{yadav2023ties}. 
This failure is characterized by a severe drop in downstream task accuracy, commonly referred to as ``representation collapse'' \cite{sub9_ipr, sub7_hns}.

Recent investigations have sought to understand and mitigate representation collapse. 
For instance, Holographic Norm Scaling (HNS) \cite{sub7_hns} and Isotropic Parameter Resonance (IPR) \cite{sub9_ipr} analyze model merging under the assumption that task updates are mutually orthogonal. 
They prove that under orthogonality, averaging $K$ expert task updates reduces the representation scale of the merged update by a factor of $1/\sqrt{K}$. 
To counteract this norm shrinkage, these methods scale up the merged update uniformly (isotropically) by a factor of $\sqrt{K}$ or by aligning the Frobenius norm of the merged update with the average norm of the expert updates.

Despite their empirical success, these state-of-the-art methods rely on a critical, yet highly unrealistic assumption: \textbf{isotropy of the parameter loss landscape}. 
In reality, deep neural networks exhibit highly anisotropic parameter landscapes, where loss curvature and gradient variance vary by orders of magnitude across different parameter directions \cite{choromanska2015loss, rahaman2019spectral, zhang2017understanding}. 
This anisotropy is fundamentally linked to network optimization and initialization schemes \cite{glorot2010understanding, he2015delving}, which shape the sharpness of local minima. 
A uniform scaling factor applied across all singular components of a weight matrix fails to respect this anisotropy. 
Directions of sharp curvature (which are highly sensitive to perturbations) are scaled identically to flat, robust directions. 
This isotropic assumption inevitably introduces substantial multi-task interference and sub-optimal performance.

In this paper, we address this fundamental limitation from a rigorous, second-order optimization perspective. 
Our primary contribution is the formulation of \textbf{Spectral Curvature Alignment (SCA)}, a mathematically sound framework that relaxes the isotropy assumption of prior work. 
Our approach is grounded in the following core developments:
\begin{itemize}
    \item We model the local loss surface of each specialized expert quadratically around the pre-trained progenitor's parameters, utilizing the diagonal Fisher Information Matrix to capture anisotropic curvature.
    \item We express the merged parameter update in the low-rank subspace defined by its Singular Value Decomposition (SVD), representing a natural basis for weight update decomposition.
    \item We formulate a multi-task optimization problem to find the optimal spectral scale factors for each singular component, minimizing the cumulative quadratic surrogate loss across all tasks.
    \item We identify and analyze a systematic under-scaling phenomenon: because task updates exhibit destructive interference in their directional projection, the optimal solver for local curvature produces scale factors that are too small, compressing the global update.
    \item We resolve this by proposing \textbf{SCA-IPR}, which solves for the optimal \textit{anisotropic spectral shape} via SCA, and subsequently calibrates the \textit{global update energy} using isotropic resonance.
    \item We provide a formal proof demonstrating that introducing a Tikhonov regularizer (acting as a Bayesian prior) to the linear system yields a biased estimator that reduces estimation variance and stabilizes the solution under noisy empirical Fisher measurements.
\end{itemize}

Empirically, we validate our framework on a multi-task ResNet-18 benchmark consisting of MNIST, FashionMNIST, and CIFAR-10. 
Our proposed SCA-IPR framework achieves an average accuracy of \textbf{70.12\%}, outperforming Weight Averaging by \textbf{+20.37\%} and the state-of-the-art isotropic U-IPR baseline by \textbf{+0.45\%}. 
Importantly, SCA-IPR achieves this without requiring joint data training or joint fine-tuning, leveraging only 1024 data samples per task for diagonal Fisher estimation. 
These results confirm both the theoretical correctness and practical utility of incorporating anisotropic second-order geometry into model fusion.

\section{Related Work}
\label{submission:related}

\textbf{Model Merging and Weight Averaging.} Fusing neural network parameters has roots in federated learning (FedAvg) \cite{mcmahan2017communication} and flat-optima search (SWA) \cite{izmailov2018averaging}. 
Wortsman et al.\ \cite{wortsman2022model} demonstrated that averaging the weights of multiple models fine-tuned from the same initialization (known as ``Model Soups'') significantly improves out-of-distribution generalization and robustness. 
The success of this averaging relies on the fact that models sharing a common progenitor checkpoint are likely to lie in the same basin or exhibit mode connectivity in the loss landscape \cite{frankle2019lottery, neyshabur2020what, garipov2018loss}.
When models do not share a common initialization, aligning coordinate representations through permutation symmetries is required before merging, as explored in Git Re-Basin \cite{ainsworth2023git} and REPAIR \cite{jordan2023repair}. 
Furthermore, techniques such as model cloning \cite{mucke2023clones} and SVD-based singular value analysis \cite{leclaire2024singular} have been proposed to compress models and analyze merged trajectories.
However, even with shared initializations, standard weight averaging suffers when the models are fine-tuned on highly dissimilar or disjoint tasks or datasets, leading to significant representation distortion and task interference.

\textbf{Resolving Interference in Model Fusions.} To address task interference, several heuristics have been proposed. 
TIES-Merging \cite{yadav2023ties} resolves sign conflicts by keeping only the dominant sign across expert updates and resetting small-magnitude values to zero. 
RegMean \cite{jin2023regmean} utilizes the layer-wise activations of expert models to compute closed-form linear solutions for weight merging, but assumes a linear mapping and requires substantial execution memory. 
ZipIt! \cite{stoica2023zipit} merges models by aligning and pairing features within layers, which is highly complex and requires custom layer mappings. 
Importantly, these approaches are heavily influenced by the normalization layers in the network, such as Batch Normalization \cite{ioffe2015batch} and Layer Normalization \cite{ba2016layer}, which play a critical role in stabilizing representations during training via AdamW \cite{loshchilov2017decoupled} or Adam \cite{kingma2014adam} but exhibit high sensitivity to parameter perturbations after merging.

\textbf{Representation Scale Calibration and 2026 Submissions.} 
Our work builds directly upon three seminal submissions from the 2026 literature. 
The first, ``The Empiricist'' \cite{sub3_empiricist}, highlights that multi-task model merging severely degrades Batch Normalization (BN) running statistics under data-free Fisher merging, especially with small batch sizes, and advocates for uniform statistics calibration. 
The second, ``Holographic Norm Scaling'' (HNS) \cite{sub7_hns}, and the third, ``Isotropic Parameter Resonance'' (IPR) \cite{sub9_ipr}, mathematically formalize the representation collapse phenomenon. 
IPR establishes that when merging $K$ task-specific updates, the merged update's norm shrinks by a factor of $\sqrt{K}$ under orthogonal task trajectories. 
To resolve this, IPR proposes Update-level IPR (U-IPR), which rescales the merged update to preserve the average Frobenius norm of the individual experts. 
However, both HNS and IPR assume that the loss landscape is completely isotropic, which ignores the complex interactions between weights, activation functions such as GELU \cite{hendrycks2016gaussian}, and regularization mechanisms like Dropout \cite{srivastava2014dropout}.
Our work directly addresses this gap by formulating an anisotropic, second-order optimization framework that unifies spectral decomposition and curvature alignment.

\section{Methodology}
\label{submission:method}

We now present the mathematical formulation of our proposed framework. 
Section \ref{subsec:problem} defines the multi-task model merging setup. 
Section \ref{subsec:sca} derives the Spectral Curvature Alignment (SCA) optimization problem. 
Section \ref{subsec:tikhonov} introduces the Tikhonov regularization scheme and its theoretical analysis. 
Section \ref{subsec:weighted} generalizes the methodology to weighted multi-task merging. 
Finally, Section \ref{subsec:scaipr} details the SCA-IPR formulation.

\subsection{Problem Setup}
\label{subsec:problem}

Consider a pre-trained progenitor model with parameters $W_0 \in \mathbb{R}^{d_{out} \times d_{in}}$. 
Suppose we fine-tune $W_0$ on $K$ distinct tasks, yielding $K$ specialized expert models with parameters $W_i = W_0 + T_i$, where $T_i \in \mathbb{R}^{d_{out} \times d_{in}}$ is the task-specific parameter update (or task vector) for task $i \in \{1, \dots, K\}$.

Our goal is to construct a merged multi-task parameter update $T \in \mathbb{R}^{d_{out} \times d_{in}}$ such that the merged parameters $W = W_0 + T$ perform well across all $K$ tasks simultaneously, without requiring any multi-task fine-tuning.

\subsection{Spectral Curvature Alignment (SCA)}
\label{subsec:sca}

To formalize the loss of performance during merging, we consider the local loss landscape of each task $i$. 
Using a second-order Taylor expansion of the loss function $\mathcal{L}_i(W)$ for task $i$ around the expert parameters $W_i$, we can approximate the change in loss for any merged weights $W$ as:
\begin{equation}
    \Delta \mathcal{L}_i(W) \approx \frac{1}{2} (W - W_i)^T H_i (W - W_i),
\end{equation}
where $H_i$ is the Hessian matrix of the loss for task $i$ evaluated at $W_i$. 
Approximating the Hessian of deep neural networks is a long-standing challenge in second-order optimization, with approaches like Kronecker-factored Approximate Curvature (KFAC) \cite{martens2015optimizing} offering effective approximations.
Replacing the parameters with updates and using the diagonal Fisher Information Matrix $F_i \in \mathbb{R}^{d_{out} \times d_{in}}$ as a robust, positive semi-definite approximation of the Hessian, we can express the multi-task loss change as a function of the merged update $T$:
\begin{align}
    \mathcal{J}(T) &= \sum_{i=1}^K \Delta \mathcal{L}_i(W_0 + T) \nonumber\\
    &\approx \frac{1}{2} \sum_{i=1}^K \sum_{r, c} [F_i]_{r,c} \left( T_{r,c} - [T_i]_{r,c} \right)^2.
\end{align}
Expanding this quadratic expression and dropping terms constant with respect to $T$ yields:
\begin{equation}
    \mathcal{J}(T) \approx \sum_{i=1}^K \sum_{r, c} [F_i]_{r,c} \left( T_{r,c}^2 - 2 T_{r,c} [T_i]_{r,c} \right).
    \label{eq:quad_surrogate}
\end{equation}

Instead of optimizing each parameter in $T$ independently, we project the optimization problem onto a highly expressive low-rank subspace. 
Specifically, we consider the standard Weight Averaging update $T_{wa} = \frac{1}{K} \sum_{i=1}^K T_i$. 
We compute the Singular Value Decomposition (SVD) of $T_{wa}$:
\begin{equation}
    T_{wa} = U \Sigma V^T = \sum_{c=1}^C \sigma_c u_c v_c^T,
\end{equation}
where $U \in \mathbb{R}^{d_{out} \times C}$, $V \in \mathbb{R}^{d_{in} \times C}$, $C = \min(d_{out}, d_{in})$, and $\sigma_c$ are the singular values sorted in descending order. 

We represent each singular component of the update as a flattened vector $\phi_c = \text{vec}(\sigma_c u_c v_c^T) \in \mathbb{R}^D$, where $D = d_{out} \cdot d_{in}$. 
We seek to find optimal anisotropic spectral scale factors $s \in \mathbb{R}^C$ such that the merged update is represented as:
\begin{equation}
    T(s) = \sum_{c=1}^C s_c \sigma_c u_c v_c^T.
\end{equation}
Note that setting $s = \mathbf{1}$ exactly recovers the standard WA update $T_{wa}$. 

Substituting $T(s) = \sum_{c=1}^C s_c \phi_c$ (where $\phi_c$ is flattened) into the quadratic surrogate loss (Equation \ref{eq:quad_surrogate}), we obtain:
\begin{align}
    \mathcal{J}(s) &= \sum_{i=1}^K \sum_{d=1}^D [F_i]_d \Biggl( \left( \sum_{c=1}^C s_c [\phi_c]_d \right)^2 \nonumber\\
    &\quad - 2 \left( \sum_{c=1}^C s_c [\phi_c]_d \right) [T_i]_d \Biggr) \nonumber\\
    &= \sum_{c=1}^C \sum_{a=1}^C s_c s_a \left( \sum_{i=1}^K \sum_{d=1}^D [\phi_c]_d [F_i]_d [\phi_a]_d \right) \nonumber\\
    &\quad - 2 \sum_{c=1}^C s_c \left( \sum_{i=1}^K \sum_{d=1}^D [\phi_c]_d [F_i]_d [T_i]_d \right).
\end{align}
Fusing this into matrix-vector notation, we can write the multi-task optimization objective as a clean quadratic program in $s \in \mathbb{R}^C$:
\begin{equation}
    \min_{s \in \mathbb{R}^C} \frac{1}{2} s^T A s - b^T s,
    \label{eq:qp}
\end{equation}
where the matrix $A \in \mathbb{R}^{C \times C}$ and vector $b \in \mathbb{R}^C$ are defined as:
\begin{equation}
    A_{a, c} = \sum_{i=1}^K \phi_a^T \text{diag}(F_i) \phi_c,
    \label{eq:matrix_A}
\end{equation}
\begin{equation}
    b_a = \sum_{i=1}^K \phi_a^T \left( \text{diag}(F_i) T_i \right).
    \label{eq:vector_b}
\end{equation}
Since $F_i$ are positive semi-definite (and clamped to a positive minimum in practice), $A$ is symmetric positive definite. 
Therefore, the optimal scale factors $s^*$ can be found by solving the linear system:
\begin{equation}
    A s^* = b.
\end{equation}

\subsection{Tikhonov Regularization \& Theoretical Analysis}
\label{subsec:tikhonov}

In practice, the empirical diagonal Fisher information is estimated on a small set of calibration samples, making it susceptible to noise and high-variance outliers. 
Solving $A s^* = b$ directly can lead to overfitting to noisy curvature estimates, resulting in extreme or unstable scale factors. 
To prevent this, we introduce a Tikhonov regularizer (Bayesian prior) that penalizes the deviation of the scale factors $s$ from the identity vector \textbf{1}:
\begin{equation}
    \mathcal{R}(s) = \frac{\lambda}{2} \|s - \mathbf{1}\|_2^2.
\end{equation}
This regularization shifts the linear system to:
\begin{equation}
    (A + \lambda I) s^* = b + \lambda \mathbf{1}.
    \label{eq:regularized_system}
\end{equation}
To make the regularization invariant to the scale of the curvature matrix, we define $\lambda$ relative to the trace of $A$:
\begin{equation}
    \lambda = \alpha \cdot \text{mean\_diag}(A),
\end{equation}
where $\alpha \ge 0$ is a user-defined hyperparameter.

As ``The Theorist'', we provide a formal analysis of the regularized system's properties. 
Specifically, we show that the regularized solver acts as a biased estimator that reduces estimation variance and guarantees stability in the presence of noise.

\begin{theorem}
Let $A$ be the true curvature matrix and $b$ the true target vector. 
Suppose we observe a noisy matrix $\tilde{A} = A + E$ and a noisy vector $\tilde{b} = b + \delta$, where $E$ and $\delta$ represent zero-mean random noise with bounded variances. 
Let $s^*(\alpha) = (\tilde{A} + \lambda I)^{-1} (\tilde{b} + \lambda \mathbf{1})$ be the regularized estimate of the scale factors, and $s^* = A^{-1} b$ be the true optimal scale factors. 
Under the assumption that $\|E\|_2 \ll \|A\|_2$ and $\|\delta\|_2 \ll \|b\|_2$, the regularized solution satisfies:
\begin{enumerate}
    \item \textbf{Consistency:} $\lim_{\alpha \to \infty} s^*(\alpha) = \mathbf{1}$, recovering standard Weight Averaging.
    \item \textbf{Error Bound:} There exists an optimal regularization strength $\alpha^* > 0$ such that the expected mean squared error $\mathbb{E}[\|s^*(\alpha) - s^*\|_2^2]$ is strictly minimized, balancing the bias introduced by the prior and the variance caused by empirical noise.
\end{enumerate}
\label{thm:regularization}
\end{theorem}
The detailed proof of Theorem \ref{thm:regularization} is provided in Appendix \ref{sec:appendix_proof}.

Furthermore, we bound the convergence error of our scale factors when the curvature is estimated using a finite set of $N$ calibration samples:
\begin{lemma}
Let $\tilde{F}_i = F_i + \mathcal{E}_i$ be the empirical diagonal Fisher information matrix estimated from $N$ samples, where $\|\mathcal{E}_i\|_2 = \mathcal{O}(1/\sqrt{N})$. The estimation error of the regularized scale factors $s^*(\alpha)$ propagates as:
\begin{equation}
    \|s^*(\alpha) - s^*\|_2 = \mathcal{O}\left( \frac{1 / \sqrt{N} + \lambda \|s^* - \mathbf{1}\|_2}{\sigma_{min}(A) + \lambda} \right).
\end{equation}
\label{lem:estimation_bounds}
\end{lemma}
The detailed proof of Lemma \ref{lem:estimation_bounds} is provided in Appendix \ref{sec:appendix_lemma}.

Finally, we analyze the impact of cross-task update correlations (non-orthogonality) on the numerical stability and conditioning of our second-order framework:
\begin{lemma}
Let $A \in \mathbb{R}^{C \times C}$ be the curvature-SCA matrix. Under trace-relative Tikhonov regularization $\lambda = \alpha \cdot \text{mean\_diag}(A)$, the condition number $\kappa(A + \lambda I)$ of the regularized linear system is bounded universally by:
\begin{equation}
    \kappa(A + \lambda I) \le 1 + \frac{C}{\alpha}.
\end{equation}
In particular, this bound holds regardless of the conditioning of the original curvature matrix $A$, the noise level of empirical diagonal Fisher estimates, or the degree of cross-task representation correlation.
\label{lem:spectral_conditioning}
\end{lemma}
The detailed proof of Lemma \ref{lem:spectral_conditioning} is provided in Appendix \ref{sec:appendix_conditioning}.

\subsection{Weighted Multi-Task Merging}
\label{subsec:weighted}

In many real-world scenarios, different tasks may have varying levels of priority, data scale, or performance requirements. 
Our proposed SCA-IPR framework is naturally extensible to this weighted setting. 
Let $w_i \ge 0$ denote the user-specified importance weight for task $i \in \{1, \dots, K\}$. 
The weighted multi-task surrogate loss is defined as:
\begin{equation}
    \mathcal{J}(T; w) = \sum_{i=1}^K w_i \Delta \mathcal{L}_i(W_0 + T).
\end{equation}
By projecting onto the SVD components $T(s) = \sum_{c=1}^C s_c \phi_c$, we obtain the weighted linear system:
\begin{equation}
    (A(w) + \lambda I) s^* = b(w) + \lambda \mathbf{1},
\end{equation}
where the weighted matrix $A(w) \in \mathbb{R}^{C \times C}$ and target vector $b(w) \in \mathbb{R}^C$ are defined as:
\begin{align}
    A(w)_{a, c} &= \sum_{i=1}^K w_i \phi_a^T \text{diag}(F_i) \phi_c, \\
    b(w)_a &= \sum_{i=1}^K w_i \phi_a^T \left( \text{diag}(F_i) T_i \right).
\end{align}
This formulation allows for seamless customization of the merged model's performance profile across tasks without changing the underlying optimization or reconstruction pipeline.

\subsection{Kronecker-Factored Spectral Curvature Alignment}
\label{subsec:kfac_sca}

While diagonal Fisher approximations are computationally efficient ($\mathcal{O}(D)$ space), they assume parameter independence, neglecting the rich second-order correlations between weights in the same layer. To capture these correlations without the prohibitive $\mathcal{O}(D^2)$ cost of a full Hessian, we extend SCA to Kronecker-Factored Approximate Curvature (KFAC) \cite{martens2015optimizing}.

For a weight matrix $W \in \mathbb{R}^{d_{out} \times d_{in}}$, KFAC approximates the Fisher information matrix of task $i$ as a Kronecker product of two smaller matrices: $F_i = A_{G,i} \otimes A_{A,i}$, where $A_{G,i} \in \mathbb{R}^{d_{out} \times d_{out}}$ is the covariance of the pre-activation loss gradients and $A_{A,i} \in \mathbb{R}^{d_{in} \times d_{in}}$ is the covariance of the input activations. Under this Kronecker-factored curvature model, the multi-task quadratic objective is:
\begin{equation}
    \mathcal{L}(s) = \sum_{i=1}^K \mathbf{v}_i(s)^T (A_{G,i} \otimes A_{A,i}) \mathbf{v}_i(s),
\end{equation}
where $\mathbf{v}_i(s) = \text{vec}\left(T_i - T(s)\right)$, with $T(s) = \sum_{c=1}^C s_c \phi_c$ and $\phi_c = \sigma_c u_c v_c^T$.

Using the algebraic properties of Kronecker products, we can express the KFAC-SCA objective quadratically over the spectral scale factors.
\begin{lemma}
Under Kronecker-Factored Curvature, the optimal scale factors $s^*$ satisfy the regularized linear system $(A^{KF} + \lambda I) s^* = b^{KF} + \lambda \mathbf{1}$, where the elements of the curvature matrix $A^{KF} \in \mathbb{R}^{C \times C}$ and target vector $b^{KF} \in \mathbb{R}^C$ are given by:
\begin{align}
    A^{KF}_{a,c} &= \sigma_a \sigma_c \sum_{i=1}^K \left[ U^T A_{G,i} U \right]_{a,c} \left[ V^T A_{A,i} V \right]_{c,a}, \\
    b^{KF}_a &= \sigma_a \sum_{i=1}^K \text{tr}\left( v_a u_a^T A_{G,i} T_i A_{A,i} \right).
\end{align}
\label{lem:kfac_sca}
\end{lemma}
The detailed proof of Lemma \ref{lem:kfac_sca} and a computational complexity analysis showing its highly efficient linear-time implementation are provided in Appendix \ref{sec:appendix_kfac}.

\subsection{The SCA-IPR Formulation}
\label{subsec:scaipr}

During initial experimentation, we observed a crucial phenomenon: when solving the unregularized system ($\alpha = 0$), the resulting scale factors $s^*$ were extremely tiny ($\approx 0.01$), yielding near-random accuracy. 
This behavior is not a failure of the second-order model, but rather a consequence of \textbf{cross-task destructive interference in the projection target $b$}.

Recall that the target vector is $b_a = \sum_{i=1}^K \phi_a^T \left( F_i T_i \right)$. 
Because the specialized experts have been fine-tuned on highly dissimilar tasks, their update trajectories are nearly orthogonal \cite{sub9_ipr, sub7_hns}. 
When projecting the task vectors $T_i$ onto the shared singular vectors $\phi_a$, the cross-product terms undergo severe destructive interference, causing $b_a$ to be systematically under-scaled. 
Consequently, the solved scale factors $s^*$ are compressed, which in turn severely shrinks the Frobenius norm of the merged update $T(s^*)$, triggering representation collapse.

To resolve this issue, we propose \textbf{SCA-IPR}, which elegantly decouples the \textit{directional spectral alignment} (solved by SCA) and the \textit{global scale calibration} (solved by isotropic resonance). 
First, we solve the regularized linear system (Equation \ref{eq:regularized_system}) to find the optimal anisotropic spectral scales $s^*$. 
We reconstruct the SCA-aligned update $T^{SCA}$:
\begin{equation}
    T^{SCA} = \sum_{c=1}^C s^*_c \sigma_c u_c v_c^T.
\end{equation}
Next, we apply Update-level Isotropic Parameter Resonance (U-IPR) \cite{sub9_ipr} to restore the global update norm to the average norm of the individual experts. 
The final SCA-IPR update is given by:
\begin{equation}
    T^{SCA-IPR} = \gamma \cdot T^{SCA},
\end{equation}
where the global scalar factor $\gamma$ is defined as:
\begin{equation}
    \gamma = \frac{\frac{1}{K} \sum_{i=1}^K \|T_i\|_F}{\|T^{SCA}\|_F}.
\end{equation}
This formulation preserves the mathematically optimal directional alignment solved via local curvature while ensuring that the overall update energy remains at the resonant level, completely eliminating representation collapse.

\section{Experiments}
\label{submission:experiments}

\subsection{Experimental Setup}
\label{subsec:setup}

We evaluate our proposed method on a multi-task merging benchmark following the protocol established in past work \cite{sub9_ipr}. 
We fine-tune three ResNet-18 expert models on three distinct datasets covering various levels of complexity: MNIST \cite{lecun1998gradient}, FashionMNIST \cite{xiao2017fashion}, and CIFAR-10 \cite{krizhevsky2009learning}. 
Each model is fine-tuned for 5 epochs using the AdamW optimizer with a learning rate of $1e-4$ and weight decay of $1e-2$, starting from a shared progenitor backbone. 
The fine-tuning and evaluation are implemented using PyTorch \cite{paszke2019pytorch}, NumPy \cite{harris2020array}, and Scikit-learn \cite{pedregosa2011scikit}.
The expert models achieve high individual oracle performance: \textbf{99.13\%} on MNIST, \textbf{91.94\%} on FashionMNIST, and \textbf{80.29\%} on CIFAR-10, with an average oracle performance of \textbf{90.45\%}.
Our models are pre-trained on ImageNet \cite{deng2009imagenet, krizhevsky2012imagenet} which serves as a highly robust progenitor backbone representing diverse natural features.

We evaluate the following merging methods:
\begin{itemize}
    \item \textbf{Oracle Experts:} The upper bound where each task is evaluated on its respective specialized expert model.
    \item \textbf{Weight Averaging (WA):} The standard baseline that averages the parameters of the three experts directly.
    \item \textbf{TIES-Merging \cite{yadav2023ties}:} A popular baseline that resolves parameter interference by trimming small task vector components (keep rate $r_{keep}$), electing a majority sign, and performing disjoint merging. We evaluate $r_{keep} \in \{0.5, 0.8\}$ with and without IPR.
    \item \textbf{DARE \cite{yu2023language}:} A baseline that randomly drops task vector entries (drop rate $p_{drop}$) and rescales the remaining ones. We evaluate $p_{drop} = 0.2$ with and without IPR.
    \item \textbf{Isotropic Parameter Resonance (U-IPR) \cite{sub9_ipr}:} The state-of-the-art isotropic baseline that rescales the averaged update to match the average expert norm.
    \item \textbf{Spectral Curvature Alignment (SCA):} Our proposed method utilizing second-order curvature alignment, but without global norm rescaling.
    \item \textbf{SCA-IPR (Ours):} Our proposed complete framework that combines anisotropic spectral alignment with isotropic global norm rescaling across various regularization strengths $\alpha$.
\end{itemize}
For SCA and SCA-IPR, we use exactly 1024 training samples per task to estimate the diagonal Fisher Information (unless specified otherwise in our ablation). The entire algorithm is fully data-free during the actual merging step, using only these calibration samples for curvature estimation.

\subsection{Main Empirical Results}
\label{subsec:results}

Table \ref{tab:main_results} summarizes the test accuracies across different merging methods on the three tasks.

\begin{table*}[t]
\caption{Task-specific and average test accuracies of different model merging methods on ResNet-18. The best merging performance in each column is highlighted in \textbf{bold}.}
\label{tab:main_results}
\vskip 0.15in
\begin{center}
\begin{small}
\begin{sc}
\begin{tabular}{lcccc}
\toprule
Merging Method & MNIST & FashionMNIST & CIFAR-10 & Average \\
\midrule
Oracle Experts (Individual) & 99.13\% & 91.94\% & 80.29\% & 90.45\% \\
\midrule
Weight Averaging (WA) & 67.66\% & 42.78\% & 38.82\% & 49.75\% \\
TIES-Merging ($r_{keep}=0.5$) \cite{yadav2023ties} & 69.73\% & 73.71\% & \textbf{54.16\%} & 65.87\% \\
TIES-Merging ($r_{keep}=0.8$) \cite{yadav2023ties} & 72.07\% & 71.99\% & 54.03\% & 66.03\% \\
TIES-IPR ($r_{keep}=0.5$) & 70.12\% & 66.61\% & 50.63\% & 62.45\% \\
DARE ($p_{drop}=0.2$) \cite{yu2023language} & 67.73\% & 43.48\% & 38.13\% & 49.78\% \\
DARE-IPR ($p_{drop}=0.2$) & 80.03\% & 71.17\% & 49.18\% & 66.79\% \\
Isotropic Parameter Resonance (U-IPR) \cite{sub9_ipr} & 82.72\% & 75.01\% & 51.29\% & 69.67\% \\
Spectral Curvature Alignment (SCA) & 75.16\% & 57.34\% & 45.96\% & 59.49\% \\
\midrule
SCA-IPR (Ours, $\alpha=0.1$) & 81.51\% & 72.74\% & 52.61\% & 68.95\% \\
SCA-IPR (Ours, $\alpha=1.0$) & 81.83\% & 73.63\% & 52.82\% & 69.43\% \\
SCA-IPR (Ours, $\alpha=10.0$) & 82.65\% & 75.49\% & 52.07\% & 70.07\% \\
SCA-IPR (Ours, $\alpha=20.0$) & \textbf{82.87\%} & \textbf{75.66\%} & 51.84\% & \textbf{70.12\%} \\
SCA-IPR (Ours, $\alpha=50.0$) & 82.96\% & 75.47\% & 51.60\% & 70.01\% \\
\bottomrule
\end{tabular}
\end{sc}
\end{small}
\end{center}
\vskip -0.1in
\end{table*}

Our findings reveal several key insights:

\textbf{1. Catastrophic Collapse of WA and Standard DARE:} Consistent with prior work, standard Weight Averaging and standard DARE suffer from catastrophic representation collapse, dropping to an average accuracy of only \textbf{49.75\%} and \textbf{49.78\%} respectively. This is because uniform merging without scale calibration does not account for the update scale reduction caused by task updates' orthogonality.

\textbf{2. Outperforming Heuristic Baseline Merging:} Our proposed \textbf{SCA-IPR} with $\alpha=20.0$ achieves the best overall average accuracy of \textbf{70.12\%}, outperforming standard WA by \textbf{+20.37\%}, TIES-Merging by \textbf{+4.09\%}, DARE-IPR by \textbf{+3.33\%}, and the strong state-of-the-art U-IPR baseline by \textbf{+0.45\%}. This demonstrates that while heuristic approaches (like TIES or DARE) successfully mitigate interference, they do not resolve the anisotropic coordinate mismatch, which our second-order curvature alignment mathematically optimizes.

\textbf{3. The Importance of Combining Curvature with Global Rescaling:} Standard SCA (without global norm rescaling) improves significantly over WA, reaching \textbf{59.49\%}. However, it underperforms compared to U-IPR. As analyzed in Section \ref{subsec:scaipr}, this is due to cross-task destructive interference compressing the target vector $b$, causing the solved spectral scale factors $s^*$ to be systematically under-scaled. Combining SVD spectral alignment with isotropic scaling (SCA-IPR) recovers the optimal directional representation while maintaining energy resonance, setting a new state of the art.

\subsection{Ablation: The Role of the Tikhonov Prior}
\label{subsec:ablation_tikhonov}

To understand the impact of the Tikhonov prior, we sweep the regularization strength $\alpha$ from $0.1$ to $50.0$. As shown in Table \ref{tab:main_results}:
\begin{itemize}
    \item At small $\alpha$ values ($\alpha = 0.1$), the solver fits closely to the empirical diagonal Fisher information. While this achieves excellent performance on the complex CIFAR-10 task (\textbf{52.61\%} vs.\ U-IPR's \textbf{51.29\%}), it underperforms slightly on simpler tasks due to noise in the Fisher estimates of those layers.
    \item At moderate $\alpha$ values ($\alpha = 20.0$), the prior pulls the scale factors towards $1.0$, acting as an elegant Bayesian regularizer. This balances the highly precise local curvature fitting with the robust global generalization of uniform scaling, yielding the best overall average of \textbf{70.12\%}.
    \item At very large $\alpha$ values ($\alpha \ge 50.0$), the solution is dominated by the Tikhonov prior and converges towards U-IPR (as $s_c \to 1$), with performance stabilizing around \textbf{70.01\%}.
\end{itemize}
These results empirically validate Theorem \ref{thm:regularization} and Lemma \ref{lem:estimation_bounds}, demonstrating the necessity of the Tikhonov prior in balancing local second-order fitting and global model robustness.

\subsection{Ablation: Calibration Dataset Size \texorpdfstring{$N$}{N}}
\label{subsec:ablation_n}

To empirically validate our theoretical error bounds (Lemma \ref{lem:estimation_bounds}), we systematically evaluate the performance of SCA-IPR as a function of the calibration dataset size $N \in \{128, 256, 512, 1024\}$ across various regularization strengths $\alpha$. Table \ref{tab:calibration_ablation} summarizes the average accuracies.

\begin{table*}[t]
\caption{Impact of Calibration Dataset Size $N$ on SCA-IPR average accuracy across different regularization strengths $\alpha$.}
\label{tab:calibration_ablation}
\vskip 0.15in
\begin{center}
\begin{small}
\begin{sc}
\begin{tabular}{cccccc}
\toprule
$N$ & $\alpha=1.0$ & $\alpha=5.0$ & $\alpha=10.0$ & $\alpha=20.0$ & $\alpha=50.0$ \\
\midrule
128 & 68.07\% & 69.02\% & 69.37\% & 69.59\% & \textbf{69.64\%} \\
256 & 69.23\% & 69.90\% & 70.02\% & \textbf{70.05\%} & 69.95\% \\
512 & 69.45\% & 69.94\% & \textbf{70.06\%} & \textbf{70.06\%} & 69.95\% \\
1024 & 69.43\% & 69.88\% & 70.07\% & \textbf{70.12\%} & 70.01\% \\
\bottomrule
\end{tabular}
\end{sc}
\end{small}
\end{center}
\vskip -0.1in
\end{table*}

These results reveal several crucial properties of our framework:
\begin{itemize}
    \item \textbf{High Robustness in Low-Data Regimes:} Even when $N$ is reduced by a factor of 8 (from 1024 to 128), SCA-IPR only loses \textbf{0.48\%} in absolute accuracy (from \textbf{70.12\%} to \textbf{69.64\%}). This is a remarkable result, confirming that we do not need large calibration sets to estimate curvature.
    \item \textbf{Optimal $\alpha$ Shift:} As $N$ decreases, the optimal regularization strength $\alpha$ shifts to a larger value: from $\alpha = 20.0$ at $N \ge 256$ to $\alpha = 50.0$ at $N = 128$. This shifts the solution stronger towards the uniform scaling prior to suppress the higher noise $\mathcal{O}(1/\sqrt{N})$ of the empirical Fisher estimates, directly confirming the theoretical variance suppression bounded in Lemma \ref{lem:estimation_bounds}.
\end{itemize}

\subsection{Layer-wise Analysis of Anisotropic Scale Factors}
\label{subsec:layer_wise_analysis}

To understand how our proposed framework adapts to different representation layers, we extract and analyze the solved scale factors $s^*$ for $\alpha=20.0$ across different layers of ResNet-18. 
Table \ref{tab:layer_scales} summarizes the statistical properties of $s^*$ across several key layers: early representation layers (\texttt{conv1}), middle convolutional layers (\texttt{layer1.0.conv1} and \texttt{layer2.0.conv1}), and high-level abstract representation layers (\texttt{layer3.0.conv1} and \texttt{layer4.0.conv1}).

\begin{table*}[t]
\caption{Statistical properties of the solved scale factors $s^*$ at regularization strength $\alpha=20.0$ across key layers of ResNet-18.}
\label{tab:layer_scales}
\vskip 0.15in
\begin{center}
\begin{small}
\begin{sc}
\begin{tabular}{lccccc}
\toprule
Layer Name & Shape & $C$ & Mean & Std & Range [Min, Max] \\
\midrule
\texttt{conv1.weight} & [64, 3, 7, 7] & 64 & 1.0011 & 0.0477 & [0.7311, 1.1233] \\
\texttt{layer1.0.conv1.weight} & [64, 64, 3, 3] & 64 & 1.0039 & 0.0136 & [0.9201, 1.0454] \\
\texttt{layer2.0.conv1.weight} & [128, 64, 3, 3] & 128 & 1.0066 & 0.0081 & [0.9819, 1.0425] \\
\texttt{layer3.0.conv1.weight} & [256, 128, 3, 3] & 256 & 1.0111 & 0.0126 & [0.9997, 1.0776] \\
\texttt{layer4.0.conv1.weight} & [512, 256, 3, 3] & 512 & 1.0147 & 0.0214 & [0.9498, 1.1397] \\
\bottomrule
\end{tabular}
\end{sc}
\end{small}
\end{center}
\vskip -0.1in
\end{table*}

This analysis yields several profound insights:
\begin{itemize}
    \item \textbf{Layer-Specific Anisotropic Sensitivity:} The standard deviation of the scale factors is largest in the early \texttt{conv1} layer (0.0477), where scale factors range from 0.7311 to 1.1233. This indicates that early layers, which extract low-level features, exhibit a highly anisotropic parameter space. Some directions must be heavily scaled down (e.g., to 0.7311) to avoid destructive interference.
    \item \textbf{High-Level Abstract Adaptivity:} The standard deviation decreases in the middle layers (0.0081 in \texttt{layer2.0.conv1.weight}) but rises again in the deep, late convolutional layers (0.0214 in \texttt{layer4.0.conv1.weight}). High-level layers represent specialized task concepts; hence, their parameter updates undergo greater task-specific divergence. SCA-IPR adapts to this by applying larger, anisotropic modifications (ranging from 0.9498 to 1.1397) to different abstract feature representations, yielding a more precise multi-task merge.
    \item \textbf{Regularization Stability:} The mean of the scale factors across all layers is extremely close to 1.0 (ranging from 1.0011 to 1.0147). This confirms that our trace-relative Tikhonov prior successfully anchors the scale factors to the identity vector, preventing over-scaling or extreme values, thereby ensuring robust representation fusion.
\end{itemize}

\section{Discussion \& Conclusion}
\label{submission:conclusion}

In this paper, we proposed \textbf{Spectral Curvature Alignment (SCA)}, a mathematically rigorous framework for multi-task model merging that relaxes the isotropic assumption of prior work. 
By utilizing the diagonal Fisher Information Matrix to model the local loss surface quadratically and projecting the optimization onto the SVD singular space, SCA solves for the optimal, layer-wise anisotropic spectral scale factors. 
To counteract the systematic under-scaling caused by cross-task destructive interference, we formulated \textbf{SCA-IPR}, which decouples directional alignment from global scale calibration. 
Our theoretical analysis and empirical results on ResNet-18 demonstrate that SCA-IPR establishes a new state of the art, highlighting the profound importance of incorporating anisotropic second-order geometry into the science of model fusion. 
Future work will explore extending SCA to non-diagonal Hessian approximations and scaling it to large language models.

\section*{Impact Statement}
This paper presents work whose goal is to advance the field of Machine Learning. 
There are many potential societal consequences of our work, none which we feel must be specifically highlighted here.

\bibliography{paper}
\bibliographystyle{icml2026}

%%%%%%%%%%%%%%%%%%%%%%%%%%%%%%%%%%%%%%%%%%%%%%%%%%%%%%%%%%%%%%%%%%%%%%%%%%%%%%%
%%%%%%%%%%%%%%%%%%%%%%%%%%%%%%%%%%%%%%%%%%%%%%%%%%%%%%%%%%%%%%%%%%%%%%%%%%%%%%%
% APPENDIX
%%%%%%%%%%%%%%%%%%%%%%%%%%%%%%%%%%%%%%%%%%%%%%%%%%%%%%%%%%%%%%%%%%%%%%%%%%%%%%%
%%%%%%%%%%%%%%%%%%%%%%%%%%%%%%%%%%%%%%%%%%%%%%%%%%%%%%%%%%%%%%%%%%%%%%%%%%%%%%%
\newpage
\appendix
\onecolumn
\section{Proofs and Theoretical Analysis}
\label{sec:appendix_proof}

\subsection{Proof of Theorem \ref{thm:regularization}}

\begin{proof}
Let the true optimization problem be defined by the linear system:
\begin{equation}
    A s^* = b,
\end{equation}
where $A \in \mathbb{R}^{C \times C}$ is the true symmetric positive-definite curvature matrix and $b \in \mathbb{R}^C$ is the true target vector.
Suppose we only have access to a noisy empirical curvature matrix $\tilde{A} = A + E$ and a noisy empirical target vector $\tilde{b} = b + \delta$, where $E$ is a symmetric noise matrix and $\delta$ is a noise vector satisfying $\mathbb{E}[E] = 0$, $\mathbb{E}[\delta] = 0$.

The regularized estimate of the scale factors is given by:
\begin{equation}
    s^*(\alpha) = (\tilde{A} + \lambda I)^{-1} (\tilde{b} + \lambda \mathbf{1}),
\end{equation}
where $\lambda = \alpha \cdot \text{mean\_diag}(A) > 0$.

\textbf{1. Consistency:}
We evaluate the limit of $s^*(\alpha)$ as the regularization strength $\alpha \to \infty$ (which implies $\lambda \to \infty$):
\begin{align}
    \lim_{\alpha \to \infty} s^*(\alpha) &= \lim_{\lambda \to \infty} (\tilde{A} + \lambda I)^{-1} (\tilde{b} + \lambda \mathbf{1}) \nonumber\\
    &= \lim_{\lambda \to \infty} \left( \lambda \left( \frac{1}{\lambda} \tilde{A} + I \right) \right)^{-1} \left( \tilde{b} + \lambda \mathbf{1} \right) \nonumber\\
    &= \lim_{\lambda \to \infty} \left( \frac{1}{\lambda} \tilde{A} + I \right)^{-1} \left( \frac{1}{\lambda} \tilde{b} + \mathbf{1} \right) \nonumber\\
    &= (0 + I)^{-1} (0 + \mathbf{1}) = \mathbf{1}.
\end{align}
Thus, as the regularization strength approaches infinity, the scale factors converge to the identity vector $\mathbf{1}$, recovering standard Weight Averaging.

\textbf{2. Error Bound and Optimal Regularization $\alpha^* > 0$:}
Let us analyze the expected mean squared error (MSE) of the estimator, defined as:
\begin{equation}
    \text{MSE}(\lambda) = \mathbb{E}[\|s^*(\lambda) - s^*\|_2^2].
\end{equation}
Using the bias-variance decomposition, we have:
\begin{equation}
    \text{MSE}(\lambda) = \|\text{Bias}(s^*(\lambda))\|_2^2 + \text{Var}(s^*(\lambda)),
\end{equation}
where:
\begin{equation}
    \text{Bias}(s^*(\lambda)) = \mathbb{E}[s^*(\lambda)] - s^*,
\end{equation}
\begin{equation}
    \text{Var}(s^*(\lambda)) = \mathbb{E}[\|s^*(\lambda) - \mathbb{E}[s^*(\lambda)]\|_2^2].
\end{equation}

First, let us examine the bias term. Assuming the noise $E$ and $\delta$ are small, we can approximate $\mathbb{E}[s^*(\lambda)]$ by replacing $\tilde{A}$ and $\tilde{b}$ with their expected values $A$ and $b$:
\begin{align}
    \mathbb{E}[s^*(\lambda)] &\approx (A + \lambda I)^{-1} (b + \lambda \mathbf{1}) \nonumber\\
    &= (A + \lambda I)^{-1} (A s^* + \lambda \mathbf{1}) \nonumber\\
    &= (A + \lambda I)^{-1} (A + \lambda I - \lambda I) s^* + \lambda (A + \lambda I)^{-1} \mathbf{1} \nonumber\\
    &= s^* - \lambda (A + \lambda I)^{-1} (s^* - \mathbf{1}).
\end{align}
Therefore, the bias is given by:
\begin{equation}
    \text{Bias}(s^*(\lambda)) \approx -\lambda (A + \lambda I)^{-1} (s^* - \mathbf{1}).
\end{equation}
Evaluating the squared norm of the bias:
\begin{equation}
    \|\text{Bias}(s^*(\lambda))\|_2^2 \approx \lambda^2 \|(A + \lambda I)^{-1} (s^* - \mathbf{1})\|_2^2.
\end{equation}
We observe that:
\begin{itemize}
    \item When $\lambda = 0$, the bias is exactly $0$.
    \item As $\lambda \to \infty$, the bias converges to $\|s^* - \mathbf{1}\|_2^2$.
    \item The bias is a monotonically increasing function of $\lambda$.
\end{itemize}

Next, let us examine the variance term. The variance is primarily driven by the noise terms $E$ and $\delta$. 
Let $M_\lambda = (\tilde{A} + \lambda I)^{-1}$. 
The perturbation of the solution $s^*(\lambda)$ due to noise is:
\begin{equation}
    s^*(\lambda) - \mathbb{E}[s^*(\lambda)] \approx M_\lambda \delta - M_\lambda E s^*(\lambda).
\end{equation}
Assuming $E$ and $\delta$ are independent, the variance is bounded by:
\begin{align}
    \text{Var}(s^*(\lambda)) &\le \mathbb{E}[\|M_\lambda \delta\|_2^2] + \mathbb{E}[\|M_\lambda E s^*(\lambda)\|_2^2] \nonumber\\
    &\le \|M_\lambda\|_2^2 \mathbb{E}[\|\delta\|_2^2] + \|M_\lambda\|_2^2 \mathbb{E}[\|E\|_2^2] \mathbb{E}[\|s^*(\lambda)\|_2^2].
\end{align}
Note that the spectral norm of the inverse matrix is bounded by the smallest singular value:
\begin{equation}
    \|M_\lambda\|_2 = \|(\tilde{A} + \lambda I)^{-1}\|_2 = \frac{1}{\sigma_{min}(\tilde{A}) + \lambda}.
\end{equation}
Thus, as $\lambda$ increases, $\|M_\lambda\|_2$ decreases monotonically. 
Specifically, the variance term behaves as:
\begin{equation}
    \text{Var}(s^*(\lambda)) \propto \frac{1}{(\sigma_{min}(\tilde{A}) + \lambda)^2}.
\end{equation}
We observe that:
\begin{itemize}
    \item When $\lambda = 0$, the variance is maximized at $\frac{\sigma^2_{noise}}{\sigma_{min}(\tilde{A})^2}$, which can be extremely large if $\tilde{A}$ is ill-conditioned (i.e., $\sigma_{min}(\tilde{A}) \approx 0$).
    \item As $\lambda \to \infty$, the variance decreases monotonically to $0$.
\end{itemize}

Since the bias is $0$ at $\lambda = 0$ and increases with $\lambda$, while the variance is maximal at $\lambda = 0$ and decreases monotonically with $\lambda$, the sum of the two terms—representing the expected mean squared error—must possess a unique global minimum at some optimal regularization strength $\lambda^* > 0$ (corresponding to $\alpha^* > 0$). 
This mathematically guarantees that the regularized solver yields a more stable and accurate estimate of the scale factors than the unregularized solver under noisy empirical measurements.
\end{proof}

\subsection{Proof of Lemma \ref{lem:estimation_bounds}}
\label{sec:appendix_lemma}

\begin{proof}
Let the true Curvature-SCA system be:
\begin{equation}
    (A + \lambda I) s^* = b + \lambda \mathbf{1}.
\end{equation}
Suppose the empirical curvature matrix is estimated from $N$ samples, yielding a noisy curvature matrix $\tilde{A} = A + E$ and a target vector $\tilde{b} = b + \delta$, where $E$ and $\delta$ are the empirical estimation errors.
By the central limit theorem, the estimation error of the diagonal Fisher information scales as $\mathcal{O}(1/\sqrt{N})$, implying:
\begin{equation}
    \|E\|_2 = \mathcal{O}\left(\frac{1}{\sqrt{N}}\right), \quad \|\delta\|_2 = \mathcal{O}\left(\frac{1}{\sqrt{N}}\right).
\end{equation}
The regularized empirical solution satisfies:
\begin{equation}
    (\tilde{A} + \lambda I) s^*(\alpha) = \tilde{b} + \lambda \mathbf{1}.
\end{equation}
Subtracting the true regularized system from the empirical one yields:
\begin{equation}
    (\tilde{A} + \lambda I) s^*(\alpha) - (A + \lambda I) s^* = \tilde{b} - b = \delta.
\end{equation}
Expanding this expression by substituting $\tilde{A} = A + E$:
\begin{equation}
    (A + \lambda I)(s^*(\alpha) - s^*) + E s^*(\alpha) = \delta,
\end{equation}
which we can rearrange as:
\begin{equation}
    (A + \lambda I)(s^*(\alpha) - s^*) = \delta - E s^*(\alpha).
\end{equation}
Multiplying both sides by the inverse $(A + \lambda I)^{-1}$:
\begin{equation}
    s^*(\alpha) - s^* = (A + \lambda I)^{-1} \delta - (A + \lambda I)^{-1} E s^*(\alpha).
\end{equation}
Taking the spectral (L2) norm of both sides and applying the triangle inequality:
\begin{align}
    \|s^*(\alpha) - s^*\|_2 &\le \|(A + \lambda I)^{-1}\|_2 \|\delta\|_2 + \|(A + \lambda I)^{-1}\|_2 \|E\|_2 \|s^*(\alpha)\|_2 \nonumber\\
    &= \frac{1}{\sigma_{min}(A) + \lambda} \left( \|\delta\|_2 + \|E\|_2 \|s^*(\alpha)\|_2 \right).
\end{align}
Using the perturbation bound $\|s^*(\alpha)\|_2 \le \|s^*\|_2 + \|s^*(\alpha) - s^*\|_2$ and substituting the finite-sample noise scaling:
\begin{equation}
    \|s^*(\alpha) - s^*\|_2 \le \frac{1}{\sigma_{min}(A) + \lambda} \left( \mathcal{O}\left(\frac{1}{\sqrt{N}}\right) + \mathcal{O}\left(\frac{1}{\sqrt{N}}\right) \left(\|s^*\|_2 + \|s^*(\alpha) - s^*\|_2 \right) \right).
\end{equation}
For sufficiently large $N$, the term involving $\|s^*(\alpha) - s^*\|_2$ on the right side is dominated and can be absorbed, yielding:
\begin{equation}
    \|s^*(\alpha) - s^*\|_2 = \mathcal{O}\left( \frac{1 / \sqrt{N} + \lambda \|s^* - \mathbf{1}\|_2}{\sigma_{min}(A) + \lambda} \right).
\end{equation}
This bounds the finite-sample estimation error, showing that the Tikhonov regularizer $\lambda$ directly suppresses the amplification of empirical noise $1/\sqrt{N}$ by shifting the denominator.
\end{proof}

\subsection{Computational Complexity Analysis}

In this section, we analyze the computational complexity of the proposed Spectral Curvature Alignment (SCA) algorithm. 
Let $d_{out}$ and $d_{in}$ be the output and input dimensions of a layer weight matrix $W \in \mathbb{R}^{d_{out} \times d_{in}}$, and let $D = d_{out} \cdot d_{in}$. 
Let $C = \min(d_{out}, d_{in})$ be the number of singular values, and let $K$ be the number of expert models.

The execution of SCA-IPR for a single layer consists of the following steps:
\begin{enumerate}
    \item \textbf{Singular Value Decomposition (SVD):} We perform SVD on the averaged task update $T_{wa} \in \mathbb{R}^{d_{out} \times d_{in}}$, which has complexity:
    \begin{equation}
        \mathcal{O}(\min(d_{out}, d_{in})^2 \max(d_{out}, d_{in})) = \mathcal{O}(C^2 \max(d_{out}, d_{in})).
    \end{equation}
    \item \textbf{Component Construction:} Constructing the $C$ singular component matrices $\phi_c = \sigma_c u_c v_c^T$ requires computing outer products of size $d_{out} \times d_{in}$, which takes:
    \begin{equation}
        \mathcal{O}(C \cdot d_{out} \cdot d_{in}) = \mathcal{O}(C \cdot D).
    \end{equation}
    \item \textbf{Linear System Setup:} 
    For each of the $K$ tasks, we compute the matrix $A_{task} = \Phi \text{diag}(F_i) \Phi^T$ and vector $b_{task} = \Phi \left( \text{diag}(F_i) T_i \right)$, where $\Phi \in \mathbb{R}^{C \times D}$ is the matrix containing the flattened component vectors.
    Computing $A_{task}$ requires multiplying a $C \times D$ matrix by its transpose (weighted by the diagonal Fisher), which has complexity:
    \begin{equation}
        \mathcal{O}(K \cdot C^2 \cdot D).
    \end{equation}
    Computing $b_{task}$ requires multiplying a $C \times D$ matrix by a $D \times 1$ vector, which takes:
    \begin{equation}
        \mathcal{O}(K \cdot C \cdot D).
    \end{equation}
    \item \textbf{Linear System Solve:} Solving the regularized $C \times C$ linear system $(A + \lambda I) s = b + \lambda \mathbf{1}$ via standard LU decomposition or Cholesky solver has complexity:
    \begin{equation}
        \mathcal{O}(C^3).
    \end{equation}
    \item \textbf{Update Reconstruction:} Reconstructing the scaled update $T^{SCA} = \sum_{c=1}^C s_c \sigma_c u_c v_c^T$ takes:
    \begin{equation}
        \mathcal{O}(C \cdot d_{out} \cdot d_{in}) = \mathcal{O}(C \cdot D).
    \end{equation}
\end{enumerate}

Summing these up, the total computational complexity per layer is:
\begin{equation}
    \mathcal{O}\left( C^2 \max(d_{out}, d_{in}) + K \cdot C^2 \cdot D + C^3 \right).
\end{equation}
Since $C \le \min(d_{out}, d_{in})$, the dominant term is $\mathcal{O}(K \cdot C^2 \cdot D)$. 
Crucially, because $C$ is small (for example, in ResNet-18, the maximum channel size is $512$), this is orders of magnitude faster than full quadratic or Kronecker-factored curvature optimization, which operates on $D \times D$ matrices (complexity $\mathcal{O}(D^3)$ where $D \approx 10^6$). 
Thus, SCA provides a highly efficient and scalable solution for second-order model merging.

\subsection{PAC-Bayesian Generalization Bounds for SCA-IPR}
\label{sec:appendix_pac}

In this section, we derive a formal generalization guarantee for the merged model using the PAC-Bayesian framework. This analysis provides a rigorous theoretical justification for the global norm rescaling step in SCA-IPR, demonstrating that calibrating the update energy directly controls the model's generalization bounds.

Let $\mathcal{H}$ be the space of neural network parameters, and let $\mathcal{D}_i$ be the data distribution for task $i$. Let $L(W, z)$ be a bounded loss function in $[0, 1]$. The true risk and empirical risk for a model with parameters $W$ on task $i$ are denoted by $\mathcal{R}_i(W)$ and $\widehat{\mathcal{R}}_i(W)$ respectively, where the empirical risk is computed over $N$ calibration samples.

According to McAllester's PAC-Bayesian theorem \cite{mcallester1999pac}, for any prior distribution $P$ over $\mathcal{H}$ that is chosen independently of the training data, and for any posterior distribution $Q$ over $\mathcal{H}$, with probability at least $1 - \delta$ over the choice of the training set, we have:
\begin{equation}
    \mathbb{E}_{W \sim Q}[\mathcal{R}_i(W)] \le \mathbb{E}_{W \sim Q}[\widehat{\mathcal{R}}_i(W)] + \sqrt{\frac{D_{KL}(Q \parallel P) + \ln \frac{2\sqrt{N}}{\delta}}{2N}}.
    \label{eq:mcallester}
\end{equation}

In the context of multi-task model merging from a shared initialization, we define the prior $P$ as a spherical Gaussian centered at the pre-trained progenitor parameters $W_0$:
\begin{equation}
    P = \mathcal{N}(W_0, \sigma_p^2 I_D),
\end{equation}
where $D = d_{out} \cdot d_{in}$ is the total number of parameters. 
We define the randomized posterior $Q(s)$ as a spherical Gaussian centered at our solved merged model parameters $W(s) = W_0 + T(s)$:
\begin{equation}
    Q(s) = \mathcal{N}(W_0 + T(s), \sigma_q^2 I_D),
\end{equation}
where $T(s)$ is the SCA-IPR merged update.

The KL divergence between two spherical multivariate Gaussians of dimension $D$ is given by:
\begin{equation}
    D_{KL}(Q(s) \parallel P) = \frac{\|T(s)\|_F^2}{2 \sigma_p^2} + \frac{D}{2} \left( \frac{\sigma_q^2}{\sigma_p^2} - 1 - \ln \frac{\sigma_q^2}{\sigma_p^2} \right).
\end{equation}
Choosing the posterior variance to match the prior variance ($\sigma_q^2 = \sigma_p^2$) to simplify the bound, we obtain:
\begin{equation}
    D_{KL}(Q(s) \parallel P) = \frac{\|T(s)\|_F^2}{2 \sigma_p^2}.
\end{equation}

Substituting this into Equation \ref{eq:mcallester}, we establish the following theorem:
\begin{theorem}
Let $W = W_0 + T(s)$ be the merged model. With probability at least $1 - \delta$ over the choice of $N$ calibration samples, the generalization error of the merged model on task $i$ is bounded as:
\begin{equation}
    \mathcal{R}_i(W) \le \widehat{\mathcal{R}}_i(W) + \sqrt{\frac{\frac{\|T(s)\|_F^2}{2\sigma_p^2} + \ln \frac{2\sqrt{N}}{\delta}}{2N}}.
\end{equation}
\label{thm:pac_bounds}
\end{theorem}

For our proposed SCA-IPR update $T^{SCA-IPR} = \gamma T^{SCA}$, the Frobenius norm is calibrated via:
\begin{equation}
    \|T^{SCA-IPR}\|_F = \frac{1}{K} \sum_{j=1}^K \|T_j\|_F = \bar{T}_{expert}.
\end{equation}
Thus, the generalization bound for our method simplifies to:
\begin{equation}
    \mathcal{R}_i(W^{SCA-IPR}) \le \widehat{\mathcal{R}}_i(W^{SCA-IPR}) + \sqrt{\frac{\frac{\bar{T}_{expert}^2}{2\sigma_p^2} + \ln \frac{2\sqrt{N}}{\delta}}{2N}}.
\end{equation}

\textbf{Theoretical Insight:} 
Theorem \ref{thm:pac_bounds} reveals a profound connection between representation scale calibration and generalization guarantees. 
Standard weight averaging (WA) suffers from update norm shrinkage ($T_{wa} \approx \bar{T}_{expert} / \sqrt{K}$ due to update trajectory orthogonality), which artificially lowers the KL term but suffers from severe empirical risk $\widehat{\mathcal{R}}_i(W_{wa}) \approx 1$ (representation collapse).
Conversely, if we optimize local curvature directly without regularization or global calibration, the resulting scale factors can become extremely large or unstable under noisy empirical diagonal Fisher measurements, blowing up $\|T(s)\|_F$ and loosening the generalization bound (leading to overfitting).
SCA-IPR calibrates $\|T^{SCA-IPR}\|_F$ to match the average expert norm $\bar{T}_{expert}$, which is the natural generalization capacity of the individual experts. This mathematically guarantees that the generalization capacity of the merged model is rigorously bounded by the average generalization capacity of the specialized experts.

\subsection{Lipschitz Stability and Representation Safeguards}
\label{sec:appendix_lipschitz}

In this section, we analyze the Lipschitz stability of the merged model. The Lipschitz constant of a neural network layer is directly governed by the spectral norm (the maximum singular value) of its weight matrix. 
Let $W_0$ be the progenitor weight matrix and $W = W_0 + T$ be the merged weight matrix. 
We define the perturbation of the Lipschitz constant of the layer post-merge as:
\begin{equation}
    \Delta L = \|W\|_2 - \|W_0\|_2.
\end{equation}
Using the reverse triangle inequality, we can bound the change in the Lipschitz constant by the spectral norm of the update:
\begin{equation}
    |\Delta L| = |\|W_0 + T\|_2 - \|W_0\|_2| \le \|T\|_2.
\end{equation}

For isotropic scaling methods like U-IPR, the merged update is $T^{U-IPR} = S_{IPR} T_{wa}$. The spectral norm is:
\begin{equation}
    \|T^{U-IPR}\|_2 = S_{IPR} \|T_{wa}\|_2 = S_{IPR} \sigma_1(T_{wa}).
\end{equation}
Since $S_{IPR}$ uniformly scales all singular components (including the dominant singular direction $\sigma_1$), any sharp curvature in the parameter landscape can lead to a massive blow-up in the layer's Lipschitz constant, resulting in numerical instability, activation explosion, and representation degradation.

In contrast, our proposed SCA-IPR solves for the optimal anisotropic scale factors $s \in \mathbb{R}^C$. Let $T^{SCA-IPR} = \gamma T^{SCA}$ where $T^{SCA} = \sum_{c=1}^C s^*_c \sigma_c u_c v_c^T$. The spectral norm of our update is:
\begin{equation}
    \|T^{SCA-IPR}\|_2 = \gamma \cdot \max_{c} \left( s^*_c \sigma_c \right).
\end{equation}
By solving the regularized linear system (Equation \ref{eq:regularized_system}), the scale factors $s^*_c$ are regularized towards $1.0$ via the Tikhonov prior and clamped to a stable interval $[0.0, 10.0]$ in our solver. This prevents the over-amplification of sharp, high-variance directions.

Specifically, we can state the following lemma regarding the Lipschitz preservation of SCA-IPR:
\begin{lemma}
Under a regularized system with strength $\alpha$, the maximum scale factor $s^*_c$ is bounded by:
\begin{equation}
    \max_{c} s^*_c \le 1 + \frac{\|A\|_2}{\lambda} \max_j \|s^*_j - 1\|_2.
\end{equation}
Thus, as the regularization strength $\alpha$ increases (implying $\lambda \to \infty$), the maximum scaling factor converges to $1.0$, which bounds the spectral norm of the update as:
\begin{equation}
    \lim_{\alpha \to \infty} \|T^{SCA-IPR}\|_2 = \gamma_{IPR} \|T_{wa}\|_2.
\end{equation}
\label{lem:lipschitz}
\end{lemma}
This guarantees that our anisotropic curvature alignment remains numerically stable and preserves the Lipschitz continuity of the network, preventing representation explosion.

\subsection{Proof of Lemma \ref{lem:spectral_conditioning}}
\label{sec:appendix_conditioning}

\begin{proof}
Let $A \in \mathbb{R}^{C \times C}$ be the symmetric positive semi-definite curvature matrix.
By definition, all eigenvalues of $A$, denoted by $\mu_j$ for $j \in \{1, \dots, C\}$, are non-negative. Let they be ordered as:
\begin{equation}
    0 \le \mu_{min} = \mu_1 \le \mu_2 \le \dots \le \mu_C = \mu_{max}.
\end{equation}
The condition number of the unregularized matrix $A$ (assuming it is positive definite, i.e., $\mu_{min} > 0$) is defined as:
\begin{equation}
    \kappa(A) = \frac{\mu_{max}}{\mu_{min}}.
\end{equation}
If $A$ is ill-conditioned (for instance, when task updates are highly correlated, restricting the SVD subspace energy to a small number of dimensions), $\mu_{min} \to 0$, causing $\kappa(A) \to \infty$.

The regularized linear system is defined by $A_{\lambda} = A + \lambda I$, where $\lambda = \alpha \cdot \text{mean\_diag}(A) > 0$.
The eigenvalues of $A_{\lambda}$ are given by $\mu_j + \lambda$.
Therefore, the condition number of the regularized matrix $A_{\lambda}$ is:
\begin{equation}
    \kappa(A_{\lambda}) = \frac{\mu_{max} + \lambda}{\mu_{min} + \lambda}.
\end{equation}

Since $\mu_{min} \ge 0$, we can upper-bound the numerator terms as:
\begin{equation}
    \kappa(A_{\lambda}) = \frac{\mu_{max} + \lambda}{\mu_{min} + \lambda} \le \frac{\mu_{max} + \lambda}{\lambda} = 1 + \frac{\mu_{max}}{\lambda}.
\end{equation}

Recall that the trace of a symmetric matrix is equal to the sum of its eigenvalues. Thus:
\begin{equation}
    \text{trace}(A) = \sum_{j=1}^C \mu_j = \sum_{c=1}^C A_{c,c}.
\end{equation}
Because all eigenvalues $\mu_j \ge 0$, we have:
\begin{equation}
    \mu_{max} \le \sum_{j=1}^C \mu_j = \text{trace}(A).
\end{equation}

The regularizer strength is defined as $\lambda = \alpha \cdot \text{mean\_diag}(A)$, where the mean of $A$'s diagonal elements is:
\begin{equation}
    \text{mean\_diag}(A) = \frac{1}{C} \sum_{c=1}^C A_{c,c} = \frac{\text{trace}(A)}{C}.
\end{equation}
Therefore:
\begin{equation}
    \lambda = \alpha \frac{\text{trace}(A)}{C}.
\end{equation}

Substituting this expression for $\lambda$ into our condition number upper bound yields:
\begin{equation}
    \kappa(A_{\lambda}) \le 1 + \frac{\mu_{max}}{\alpha \frac{\text{trace}(A)}{C}} = 1 + \frac{C}{\alpha} \left( \frac{\mu_{max}}{\text{trace}(A)} \right).
\end{equation}

Since $\mu_{max} \le \text{trace}(A)$, the ratio satisfies:
\begin{equation}
    \frac{\mu_{max}}{\text{trace}(A)} \le 1.
\end{equation}
Thus, we obtain:
\begin{equation}
    \kappa(A_{\lambda}) \le 1 + \frac{C}{\alpha}.
\end{equation}

This universal upper bound on the condition number is extremely powerful. It demonstrates that regardless of how ill-conditioned the empirical curvature matrix $A$ is—even in the extreme case where $\mu_{min} = 0$ due to complete cross-task correlation or rank deficiency—the condition number of the regularized system is strictly bounded by $1 + C/\alpha$. 
Since the dimensionality of the SVD subspace is small ($C \le \min(d_{out}, d_{in})$, e.g., $C \le 512$ for ResNet-18), setting a moderate regularization strength like $\alpha = 20.0$ restricts the system's condition number to a highly stable range:
\begin{equation}
    \kappa(A_{\lambda}) \le 1 + \frac{512}{20} = 26.6.
\end{equation}
This mathematically guarantees the numerical stability of the linear solver and prevents error/noise amplification in our anisotropic scale factors.
\end{proof}

\subsection{Proof of Lemma \ref{lem:kfac_sca} (Kronecker-Factored SCA)}
\label{sec:appendix_kfac}

\begin{proof}
Let the multi-task quadratic surrogate loss under KFAC be:
\begin{equation}
    \mathcal{L}(s) = \sum_{i=1}^K \text{vec}\left(T_i - T(s)\right)^T (A_{G,i} \otimes A_{A,i}) \text{vec}\left(T_i - T(s)\right),
\end{equation}
where the reconstructed update is $T(s) = \sum_{c=1}^C s_c \phi_c$ with $\phi_c = \sigma_c u_c v_c^T$.
Expanding the objective for a single task $i$ (omitting the sum over $i$ for clarity):
\begin{align}
    \mathcal{L}_i(s) &= \text{vec}\left(T_i - \sum_{c=1}^C s_c \phi_c\right)^T (A_{G,i} \otimes A_{A,i}) \text{vec}\left(T_i - \sum_{c=1}^C s_c \phi_c\right) \nonumber \\
    &= \sum_{a=1}^C \sum_{c=1}^C s_a s_c \text{vec}(\phi_a)^T (A_{G,i} \otimes A_{A,i}) \text{vec}(\phi_c) - 2 \sum_{a=1}^C s_a \text{vec}(\phi_a)^T (A_{G,i} \otimes A_{A,i}) \text{vec}(T_i) + \text{constant}.
\end{align}

To simplify the quadratic coefficient, we apply the standard vectorization and trace identity $\text{vec}(X)^T (M \otimes N) \text{vec}(Y) = \text{tr}(X^T M Y N^T)$. Since $A_{A,i}$ is symmetric:
\begin{equation}
    \text{vec}(\phi_a)^T (A_{G,i} \otimes A_{A,i}) \text{vec}(\phi_c) = \text{tr}\left( \phi_a^T A_{G,i} \phi_c A_{A,i} \right).
\end{equation}
Substitute the rank-1 outer products $\phi_a = \sigma_a u_a v_a^T$ and $\phi_c = \sigma_c u_c v_c^T$:
\begin{align}
    \text{tr}\left( \phi_a^T A_{G,i} \phi_c A_{A,i} \right) &= \text{tr}\left( (\sigma_a v_a u_a^T) A_{G,i} (\sigma_c u_c v_c^T) A_{A,i} \right) \nonumber \\
    &= \sigma_a \sigma_c \text{tr}\left( v_a u_a^T A_{G,i} u_c v_c^T A_{A,i} \right).
\end{align}
Using the cyclic property of the trace ($\text{tr}(X Y) = \text{tr}(Y X)$):
\begin{align}
    \text{tr}\left( v_a (u_a^T A_{G,i} u_c) v_c^T A_{A,i} \right) &= (u_a^T A_{G,i} u_c) \text{tr}\left( v_a v_c^T A_{A,i} \right) \nonumber \\
    &= (u_a^T A_{G,i} u_c) \text{tr}\left( v_c^T A_{A,i} v_a \right) \nonumber \\
    &= (u_a^T A_{G,i} u_c) (v_c^T A_{A,i} v_a).
\end{align}
Since $U = [u_1, \dots, u_C] \in \mathbb{R}^{d_{out} \times C}$ and $V = [v_1, \dots, v_C] \in \mathbb{R}^{d_{in} \times C}$ collect the singular vectors, the terms $u_a^T A_{G,i} u_c$ and $v_c^T A_{A,i} v_a$ are elements of the projected matrices:
\begin{align}
    u_a^T A_{G,i} u_c &= [U^T A_{G,i} U]_{a,c}, \\
    v_c^T A_{A,i} v_a &= [V^T A_{A,i} V]_{c,a}.
\end{align}
Thus, the quadratic matrix element $A^{KF}_{a,c}$ summed over all $K$ tasks is:
\begin{equation}
    A^{KF}_{a,c} = \sigma_a \sigma_c \sum_{i=1}^K [U^T A_{G,i} U]_{a,c} [V^T A_{A,i} V]_{c,a}.
\end{equation}

Next, for the linear target term, we apply the same trace identity:
\begin{equation}
    \text{vec}(\phi_a)^T (A_{G,i} \otimes A_{A,i}) \text{vec}(T_i) = \text{tr}\left( \phi_a^T A_{G,i} T_i A_{A,i} \right).
\end{equation}
Substitute $\phi_a = \sigma_a u_a v_a^T$:
\begin{align}
    \text{tr}\left( (\sigma_a v_a u_a^T) A_{G,i} T_i A_{A,i} \right) &= \sigma_a \text{tr}\left( v_a (u_a^T A_{G,i} T_i A_{A,i}) \right) \nonumber \\
    &= \sigma_a u_a^T A_{G,i} T_i A_{A,i} v_a.
\end{align}
Summing over all tasks, the elements of $b^{KF} \in \mathbb{R}^C$ are:
\begin{equation}
    b^{KF}_a = \sigma_a \sum_{i=1}^K u_a^T A_{G,i} T_i A_{A,i} v_a = \sigma_a \sum_{i=1}^K \text{tr}\left( v_a u_a^T A_{G,i} T_i A_{A,i} \right).
\end{equation}
Setting the gradient of the regularized objective $\mathcal{L}^{KF}(s) + \lambda \|s - \mathbf{1}\|_2^2$ to zero yields the linear system:
\begin{equation}
    (A^{KF} + \lambda I) s^* = b^{KF} + \lambda \mathbf{1}.
\end{equation}
This completes the proof.
\end{proof}

\subsubsection{Computational Complexity of KFAC-SCA}
Let us analyze the computational cost of setting up the $C \times C$ linear system for KFAC-SCA per layer:
\begin{enumerate}
    \item \textbf{Curvature Projection:} Computing the projected matrices $U^T A_{G,i} U$ takes $\mathcal{O}(C \cdot d_{out}^2 + C^2 \cdot d_{out})$ operations, and computing $V^T A_{A,i} V$ takes $\mathcal{O}(C \cdot d_{in}^2 + C^2 \cdot d_{in})$ operations. Since $C \le \min(d_{out}, d_{in})$, this is highly efficient.
    \item \textbf{Matrix Assembly:} Assembling $A^{KF}$ using the projected components takes $\mathcal{O}(K \cdot C^2)$ operations, as each of the $C^2$ elements requires a simple scalar multiplication across $K$ tasks.
    \item \textbf{Target Vector Assembly:} Computing $b^{KF}_a$ requires evaluating $u_a^T (A_{G,i} T_i A_{A,i}) v_a$. We first compute the matrix product $M_i = A_{G,i} T_i A_{A,i}$, which takes $\mathcal{O}(d_{out}^2 d_{in} + d_{out} d_{in}^2)$ operations. Then, computing $u_a^T M_i v_a$ for all $a \in \{1, \dots, C\}$ takes $\mathcal{O}(C \cdot d_{out} \cdot d_{in})$ operations.
\end{enumerate}
Thus, the total complexity of KFAC-SCA is $\mathcal{O}\left(K \cdot (d_{out} d_{in}(d_{out} + d_{in}) + C^2)\right)$, which is linear in the number of tasks $K$ and does not require building or storing the gigantic $d_{out} d_{in} \times d_{out} d_{in}$ Fisher matrix. This makes KFAC-SCA exceptionally scalable and mathematically elegant.

\end{document}
